\section{Описание}

Требуется реализовать алгоритм Кнута"---Морриса"---Пратта для поиска одного образца
в тексте над алфавитом чисел от~0 до~$2^{32}-1$.

Алгоритм КМП работает за $O(|P| + |T|)$, где $|P|$~--- длина образца, а $|T|$~---
длина текста. Сначала для образца строится префикс-функция $\pi$: значение
$\pi[i]$ равно длине наибольшего собственного префикса $P[0..i]$, который одновременно
является суффиксом этой последовательности. Затем текст просматривается слева направо.
При несовпадении алгоритм не возвращается назад по тексту, а использует $\pi$ для
выбора следующей длины уже совпавшего префикса.

Числа считываются в тип \texttt{uint32\_t}. Стандартное числовое считывание корректно
обрабатывает ведущие нули: \texttt{0011} и \texttt{11} становятся одним и тем же
значением~11.

\pagebreak

\section{Исходный код}

\subsection{Префикс-функция}

Функция \texttt{BuildPrefixFunction} строит массив $\pi$ для образца. Переменная
\texttt{k} хранит текущую длину совпавшего префикса. Пока следующий символ не подходит,
значение \texttt{k} откатывается по уже посчитанным значениям префикс-функции.

\begin{lstlisting}[language=C++, caption={Построение префикс-функции}]
std::vector<int> BuildPrefixFunction(const std::vector<TToken> &pattern) {
    int patternLen = static_cast<int>(pattern.size());
    std::vector<int> prefixFunc(patternLen, 0);

    for (int i = 1; i < patternLen; ++i) {
        int k = prefixFunc[i - 1];
        while (k > 0 && pattern[i] != pattern[k]) {
            k = prefixFunc[k - 1];
        }
        if (pattern[i] == pattern[k]) {
            ++k;
        }
        prefixFunc[i] = k;
    }

    return prefixFunc;
}
\end{lstlisting}

\subsection{Поиск}

В основном решении текст обрабатывается потоково, построчно и пословно. Это важно,
поскольку в условии нет ограничений на длину текста, а лимит памяти составляет 64 MB.
Для вывода позиции найденного вхождения хранится только очередь последних $|P|$
позиций \texttt{positionWindow}. Сам алгоритм остаётся классическим КМП: переменная
\texttt{matched} хранит длину совпавшего префикса образца.

\begin{lstlisting}[language=C++, caption={Потоковый проход КМП в main}]
while (ss >> token) {
    positionWindow.emplace_back(lineNumber, wordNumber);
    if (static_cast<int>(positionWindow.size()) > patternLen) {
        positionWindow.pop_front();
    }

    while (matched > 0 && token != pattern[matched]) {
        matched = prefixFunc[matched - 1];
    }
    if (token == pattern[matched]) {
        ++matched;
    }
    if (matched == patternLen) {
        const TPosition &match = positionWindow.front();
        std::cout << match.first << ", " << match.second << "\n";
        matched = prefixFunc[matched - 1];
    }

    ++wordNumber;
}
\end{lstlisting}

Для теста производительности также реализована функция \texttt{KmpSearch}, принимающая
готовые векторы текста и позиций. Она использует тот же цикл КМП и нужна только для
удобного сравнения с наивным алгоритмом в бенчмарке.

\pagebreak

\section{Консоль}

\begin{alltt}
# cmake -B build -DCMAKE_BUILD_TYPE=Release
# cmake --build build
# ./build/lab4/lab4_main
11 45 11 45 90
0011 45 011 0045 11 45 90    11
45 11 45 90
1, 3
1, 8
\end{alltt}

\pagebreak
